%======================= RÉALISATION PAR MODULE =============================
\chapter{Réalisation technique par module}

Ce chapitre présente, module par module, l'objectif poursuivi, les
fonctionnalités implémentées, les principaux choix techniques, une capture
d'écran illustrative et, le cas échéant, les limites à ce stade.

%---------------------------------------------------------------------------
\section{Module A --- Connexion et navigation}

\subsection*{Objectif}
Offrir un point d'entrée sécurisé (dans l'esprit) à l'application et adapter la
navigation au rôle de l'utilisateur.

\subsection*{Fonctionnalités}
\begin{itemize}
  \item écran de connexion avec e-mail, mot de passe et \textbf{choix du
        rôle} (Admin, Media buyer, Lecteur) ;
  \item \textbf{navigation latérale} filtrée selon le rôle ;
  \item \textbf{protection des routes} : redirection vers la connexion si
        aucune session, ou vers le tableau de bord si le rôle n'est pas
        autorisé.
\end{itemize}

\subsection*{Explication technique}
La session est gérée par un contexte d'authentification et persistée dans le
\texttt{localStorage} du navigateur. Deux gardes de route encadrent l'accès aux
pages. Il s'agit d'une \textbf{authentification simulée} : aucun mot de passe
n'est vérifié et le contrôle d'accès n'est appliqué que côté frontend
(voir limites).

\begin{figure}[H]
\centering
\includegraphics[width=0.92\textwidth]{login.png}
\caption{Interface de connexion}
\label{fig:login}
\end{figure}

%---------------------------------------------------------------------------
\section{Module B --- Tableau de bord}

\subsection*{Objectif}
Donner une vue de synthèse immédiate de l'activité publicitaire.

\subsection*{Fonctionnalités}
\begin{itemize}
  \item quatre \textbf{indicateurs clés} (impressions, clics, budget dépensé,
        campagnes actives) ;
  \item \textbf{graphique de performance} ;
  \item \textbf{tableau des campagnes récentes}.
\end{itemize}

\subsection*{Explication technique}
Les données proviennent d'un endpoint dédié (\texttt{GET
/api/statistics/dashboard}) qui agrège des données \emph{mockées}, prêtes à
être remplacées par de vraies statistiques une fois la base branchée. Les
graphiques utilisent la bibliothèque Recharts.

\begin{figure}[H]
\centering
\includegraphics[width=0.92\textwidth]{dashboard.png}
\caption{Tableau de bord principal}
\label{fig:dashboard}
\end{figure}

%---------------------------------------------------------------------------
\section{Module C --- Gestion des campagnes}

\subsection*{Objectif}
Lister les campagnes existantes et les brouillons en cours.

\subsection*{Fonctionnalités}
\begin{itemize}
  \item section \textbf{\og Vos brouillons \fg{}} alimentée par l'API ;
  \item section \textbf{\og Toutes les campagnes \fg{}} présentant des
        campagnes de démonstration (active, en pause, brouillon, terminée) ;
  \item accès direct à l'assistant de création.
\end{itemize}

\subsection*{Explication technique et limite}
La liste des brouillons est \textbf{réellement branchée} au backend
(\texttt{GET /api/campaigns/drafts}). En revanche, la table \og Toutes les
campagnes \fg{} affiche encore des données \emph{mockées} côté frontend : ce
point est signalé en toute transparence et figure parmi les évolutions
possibles.

\begin{figure}[H]
\centering
\includegraphics[width=0.92\textwidth]{campagnes.png}
\caption{Liste des campagnes et des brouillons}
\label{fig:campagnes}
\end{figure}

%---------------------------------------------------------------------------
\section{Module D --- Assistant de création de campagne}

\subsection*{Objectif}
Guider l'utilisateur dans la création d'une campagne à travers un parcours
structuré.

\subsection*{Fonctionnalités}
Assistant en \textbf{cinq étapes} : (1) informations générales (nom,
annonceur, objectif, dates, budgets), (2) ciblage technique (appareils,
systèmes d'exploitation, plages horaires), (3) formats publicitaires, (4)
magasins ciblés (voir modules E et F), (5) catégories d'applications et
résumé. À toute étape, l'utilisateur peut \textbf{enregistrer un brouillon},
même incomplet.

\subsection*{Explication technique}
L'assistant s'appuie sur un catalogue d'options servi par l'API
(\texttt{GET /api/campaign-creation/options}), garantissant un vocabulaire
partagé entre frontend et backend. L'enregistrement d'un brouillon
(\texttt{POST /api/campaigns/drafts}) constitue la \textbf{seule donnée
réellement persistée} du prototype ; en cas d'API injoignable, un repli local
(navigateur) évite de bloquer le parcours.

\begin{figure}[H]
\centering
\includegraphics[width=0.92\textwidth]{creation-campagne.png}
\caption{Assistant de création de campagne}
\label{fig:creation}
\end{figure}

%---------------------------------------------------------------------------
\section{Module E --- Import des magasins et validation CSV}

\subsection*{Objectif}
Permettre à l'annonceur d'importer sa base de magasins et en garantir la
qualité avant utilisation.

\subsection*{Fonctionnalités}
\begin{itemize}
  \item import d'un fichier \texttt{.xlsx} ou \texttt{.csv} ;
  \item \textbf{validation ligne par ligne} : colonnes obligatoires,
        coordonnées GPS, format d'URL, détection des doublons ;
  \item affichage du nombre de lignes valides et des erreurs détaillées ;
  \item sélection des magasins à cibler.
\end{itemize}

\subsection*{Explication technique}
La validation est réalisée \textbf{côté serveur}
(\texttt{POST /api/stores/import/preview}) : le fichier est analysé (CSV avec
séparateur \texttt{,} ou \texttt{;}, ou Excel), chaque ligne est contrôlée, et
une erreur explicite est produite par cellule invalide, avec son numéro de
ligne. Aucune donnée n'est persistée à l'étape d'aperçu.

\begin{figure}[H]
\centering
\includegraphics[width=0.92\textwidth]{magasins.png}
\caption{Import et validation de la base de magasins}
\label{fig:magasins}
\end{figure}

%---------------------------------------------------------------------------
\section{Module F --- Géociblage / geofencing}

\subsection*{Objectif}
Définir et visualiser les zones de diffusion autour des magasins sélectionnés.

\subsection*{Fonctionnalités}
\begin{itemize}
  \item carte interactive avec un \textbf{marqueur par magasin} ;
  \item réglage du \textbf{rayon de diffusion} (1 à 20 km, 5 km par défaut) ;
  \item visualisation du rayon par un \textbf{cercle} sur la carte.
\end{itemize}

\subsection*{Explication technique}
La cartographie repose sur \textbf{Leaflet} et les fonds de carte
\textbf{OpenStreetMap} (aucune clé propriétaire requise). Le composant de carte
est \textbf{partagé} entre la page \og Magasins \fg{} et l'étape 4 de
l'assistant, garantissant un comportement identique aux deux endroits.

\begin{figure}[H]
\centering
\includegraphics[width=0.92\textwidth]{map-geofencing.png}
\caption{Carte de géociblage (geofencing)}
\label{fig:geofencing}
\end{figure}

%---------------------------------------------------------------------------
\section{Module G --- DCO et landing pages simulées}

\subsection*{Objectif}
Générer automatiquement des créatives et des pages d'atterrissage
personnalisées par magasin.

\subsection*{Fonctionnalités}
\begin{itemize}
  \item upload d'un visuel par format (bannière, pavé, interstitiel) ;
  \item \textbf{génération d'une variante par magasin} intégrant les champs
        dynamiques (nom, ville, adresse, horaires) ;
  \item comparaison de deux magasins et galerie complète des variantes ;
  \item \textbf{\emph{landing pages} simulées} prévisualisables par magasin.
\end{itemize}

\subsection*{Explication technique et limites}
Seules les \textbf{métadonnées} du visuel sont envoyées et enregistrées côté
serveur (\texttt{POST /api/dco/creatives}) : aucun fichier n'est stocké sur
disque. La génération des variantes et des \emph{landing pages} est calculée
\textbf{côté frontend}. Les \emph{landing pages} sont \textbf{entièrement
simulées} (pas de route publique réelle). Enfin, les URL de magasin issues des
données de démonstration sont \textbf{volontairement non cliquables} (texte
informatif), car elles ne pointent pas vers de vraies pages en ligne.

\begin{figure}[H]
\centering
\includegraphics[width=0.92\textwidth]{dco.png}
\caption{Module DCO --- génération de variantes}
\label{fig:dco}
\end{figure}

%---------------------------------------------------------------------------
\section{Module H --- Reporting et export CSV}

\subsection*{Objectif}
Analyser la performance des campagnes selon différents filtres et exporter les
résultats.

\subsection*{Fonctionnalités}
\begin{itemize}
  \item filtres (période, campagne, ville/magasin) ;
  \item \textbf{indicateurs} et \textbf{graphiques} qui se recalculent selon
        les filtres ;
  \item \textbf{carte des zones de diffusion} et \textbf{tableau par magasin}
        triable ;
  \item \textbf{export CSV} compatible Excel.
\end{itemize}

\subsection*{Explication technique et limite}
La liste des magasins et des villes provient de l'endpoint réel
(\texttt{GET /api/stores}). En revanche, les indicateurs de performance sont
\textbf{générés côté frontend} (données de démonstration réalistes), et non
encore agrégés depuis une base. L'export CSV inclut un \textbf{BOM UTF-8} et un
séparateur point-virgule ({\NoAutoSpacing\texttt{;}}), afin de garantir un
affichage correct des accents et des colonnes dans Excel.

\begin{figure}[H]
\centering
\includegraphics[width=0.92\textwidth]{reporting.png}
\caption{Reporting des performances}
\label{fig:reporting}
\end{figure}

%---------------------------------------------------------------------------
\section{Module I --- Gestion du compte}

\subsection*{Objectif}
Administrer les annonceurs, les utilisateurs et les paramètres du compte.

\subsection*{Fonctionnalités}
\begin{itemize}
  \item onglet \textbf{Annonceurs} : recherche, filtre, fiche détaillée,
        modification ;
  \item onglet \textbf{Utilisateurs} : recherche, filtre par rôle, ajout,
        modification, activation/désactivation ;
  \item onglet \textbf{Paramètres} : devise, fuseau horaire, langue, e-mail de
        notification, activation du suivi.
\end{itemize}

\subsection*{Explication technique}
Le module est \textbf{entièrement branché sur l'API}
(\texttt{/api/advertisers}, \texttt{/api/users}, \texttt{/api/account-settings}),
avec des erreurs HTTP explicites (404 pour un identifiant inconnu, 409 pour un
e-mail déjà utilisé, 422 pour un champ invalide). L'\textbf{édition} est
réservée au rôle Admin ; les autres rôles accèdent en lecture seule.

\begin{figure}[H]
\centering
\includegraphics[width=0.92\textwidth]{gestion-compte.png}
\caption{Gestion du compte (annonceurs, utilisateurs, paramètres)}
\label{fig:compte}
\end{figure}

%---------------------------------------------------------------------------
\section{Module J --- Backend FastAPI et Swagger}

\subsection*{Objectif}
Exposer une API REST claire, validée et documentée.

\subsection*{Fonctionnalités}
\begin{itemize}
  \item \textbf{24 endpoints} REST regroupés par ressource ;
  \item validation des données par Pydantic et erreurs HTTP explicites ;
  \item documentation \textbf{Swagger/OpenAPI} interactive et testable.
\end{itemize}

\subsection*{Explication technique}
FastAPI génère automatiquement la documentation OpenAPI, accessible à
l'adresse \texttt{/docs}. Chaque endpoint y est décrit et peut être exécuté
directement depuis l'interface. La figure~\ref{fig:swagger} présente la vue
d'ensemble de l'API, et la figure~\ref{fig:swaggerhealth} illustre le test de
l'endpoint de santé.

\begin{figure}[H]
\centering
\includegraphics[width=0.92\textwidth]{swagger-api.png}
\caption{Documentation Swagger de l'API}
\label{fig:swagger}
\end{figure}

\begin{figure}[H]
\centering
\includegraphics[width=0.92\textwidth]{swagger-health-test.png}
\caption{Test de l'endpoint \texttt{/api/health} depuis Swagger}
\label{fig:swaggerhealth}
\end{figure}

%---------------------------------------------------------------------------
\section{Module K --- GitHub et documentation}

\subsection*{Objectif}
Assurer un versionnement rigoureux et une documentation exploitable du projet.

\subsection*{Fonctionnalités}
\begin{itemize}
  \item dépôt GitHub structuré (\texttt{frontend/}, \texttt{backend/},
        \texttt{database/}, \texttt{docs/}) ;
  \item flux par \textbf{branches} et \textbf{\emph{pull requests}} ;
  \item documentation continue (technique, schéma de données, scénario et
        checklist de démonstration, journal d'avancement, livraison finale).
\end{itemize}

\subsection*{Explication technique}
Chaque incrément a été développé sur une branche dédiée puis intégré par
\emph{pull request}, ce qui conserve un historique lisible et réversible. La
figure~\ref{fig:github} présente le dépôt du projet.

\begin{figure}[H]
\centering
\includegraphics[width=0.92\textwidth]{github-repository.png}
\caption{Dépôt GitHub du projet}
\label{fig:github}
\end{figure}
